\section{Experimental details and reproduction index}\label{app:experimental-details}

Table~\ref{tab:reproduction-index} locates the settings and code for each study. Paths are relative to the artifact's \texttt{experiments/} directory. The HMM studies reuse one trained checkpoint; the synthetic studies evaluate an explicit response function. Repeated masks, prompt batches, or design seeds are therefore not independently trained models. Numerical precision is documented in Appendix~\ref{app:precision}.

\begin{table*}[t]
\centering
\begin{tabular}{@{}>{\raggedright\arraybackslash}p{0.30\linewidth}>{\raggedright\arraybackslash}p{0.66\linewidth}@{}}
\toprule
Study and further details & Artifact entry point \\
\midrule
HMM control (\S\ref{sec:result-control}, Appendix~\ref{app:hmm-control-setup}) & \path{hmm/hmm_observer_control.py}; \path{hmm/control_gate_d_rotating.py} \\
\addlinespace
HMM recovery (\S\ref{sec:aggregate}, Appendix~\ref{app:hmm-measurements}) & \path{hmm/nt_mi_correspondence.py}; \path{hmm/sparse_tomography_posthoc.py} \\
\addlinespace
HMM calibration (\S\ref{sec:step0}, Appendix~\ref{app:hmm-measurements}) & \path{hmm/attribution_vs_finite_step0.py}; \path{revision_checks/calibration_seed_factorial.py}; \path{hmm_selector/} \\
\addlinespace
Gemma (Appendix~\ref{app:gemma}) & \path{gemma/packet/collect.py}; \path{gemma/packet/analysis.py}; \path{gemma/reproduce.py} \\
\addlinespace
Synthetic (\S\ref{sec:interactions}, Appendix~\ref{app:synthetic-setup}) & \path{planted_interactions/claim3_planted_reach.py}; \path{hvp_interactions/claim3_hvp_baseline.py} \\
\addlinespace
Tracr (\S\ref{sec:basis}) & \path{tracr/tracr_label_basis_analysis.py}; \path{tracr/tracr_residual_group_intervention_adapter.py} \\
\addlinespace
GPT-2 IOI (\S\ref{sec:ioi}) & \path{ioi/stage2c/ioi_stage2c_primary_stratified.py}; \path{ioi/stage2d/ioi_stage2d_per_pair_decomposition.py} \\
\addlinespace
Qwen (\S\ref{sec:qwen}, Appendix~\ref{app:qwen-reproduction}) & \path{qwen/configs/qwen2_5_7b_h200_full.json}; \path{qwen/src/mechtomo/analysis.py} \\
\bottomrule
\end{tabular}
\caption{Code entry points. Main-text sections give the experimental setup; appendices provide additional settings and records.}\label{tab:reproduction-index}
\end{table*}

\subsection{HMM data, transformer, and one-step control}\label{app:hmm-control-setup}

\paragraph{Data and training.}
The binary HMMs in Sec.~\ref{sec:result-control} start from uniform state distributions. For stay probability $p_s$ and correct-emission probability $p_e$, their transition and emission matrices are
\[
T=\begin{pmatrix}p_s&1-p_s\\1-p_s&p_s\end{pmatrix},\qquad
E=\begin{pmatrix}p_e&1-p_e\\1-p_e&p_e\end{pmatrix}.
\]
A token encodes the two observations as $o_{1,t}+2o_{2,t}$.

In addition to the architecture and training settings in Sec.~\ref{sec:result-control}, \path{hmm/frozen/config.json} specifies GELU feed-forward blocks, learned positional embeddings, pre-layer normalization, and no dropout. AdamW uses weight decay $0.01$ and gradient norm clipping at 1. Training predicts tokens 2--64 from tokens 1--63 with next-token cross-entropy and no posterior supervision. Evaluation at step 1 and every 100 steps draws five further batches; the script preserves this random-number sequence. The final checkpoint is \path{hmm/frozen/model.pt}.

\paragraph{Readout and units.}
Let $q_t$ be the model's probability that the next observation from a given chain is one. If $p_t$ is the current hidden-state probability, the HMM predicts
\[
q_t=(1-p_e)+(2p_e-1)\big[(1-p_s)+(2p_s-1)p_t\big].
\]
We invert this affine relation, clip $p_t$ to $[10^{-5},1-10^{-5}]$, and take $\log(p_t/(1-p_t))$ to obtain the output-implied log-odds. The model's four-way softmax gives $q_{1,t}$ by summing token probabilities 1 and 3, and $q_{2,t}$ by summing probabilities 2 and 3. This construction explains both the readout's units and its nonsmooth clipping boundaries.

\paragraph{Trained readout and edit direction.}
For each layer, ridge regression predicts $z_1$ and $z_2$ from standardized residual activations, with penalty $0.001$ and an unpenalized intercept. Training activations come from 40 batches of 256 sequences, subsampled to 250,000 positions; validation uses ten fresh batches and 50,000 positions. The smallest validation RMSE for $z_1$ selects layer 3 (zero-based), the final block. The high-minus-low direction $d_1$ defined in Sec.~\ref{sec:result-control} is distinct from these regression weights.

If $b$ is the standardized readout's weight vector and $s_h$ its training standard deviations, the controller's gain is $g_1=\sum_i b_i d_{1,i}/s_{h,i}$, retained as $0.41545859$. The second pass adds $u_t d_1$ before the final layer normalization and output head, as specified in Sec.~\ref{sec:result-control}.

\paragraph{Metrics and repeats.}
Observer RMSE is computed against the exact posterior log-odds within each batch and then averaged. The correlation in Figure~\ref{fig:observer} uses ten observer-level means with the target and off-target metrics defined in Sec.~\ref{sec:result-control}. The observers share one trained model and seed. Their retained rows are in \path{hmm/frozen/observer_control.csv}.

\paragraph{Separate mixed-direction comparison.}
The specificity script constructs $d_2$ from high-minus-low $z_2$ activations, removes its projection onto $d_1$, and normalizes it. At $\alpha=0.5$, it uses the estimate $z_1+\alpha z_2$ and direction $(d_1+\alpha d_2)/\|d_1+\alpha d_2\|$. Its controller denominator combines the two trained readout gains with weights $1$ and $\alpha$. The recovered original CSV, \path{hmm/frozen/gate_d_rotating_v1/gate_d_rotating_control.csv}, verifies the rounded values in Sec.~\ref{sec:result-control}: mixed-case target MSE $25.0173\to4.8124$, oracle MSE $3.5623$, and off-target probability movement $0.03674\to0.07923$.

A separate CPU verification retains the checkpoint and readouts, using seed 7, 20 direction batches with at most 120,000 positions, and 20 control batches of 256 sequences per variant. It gives mixed-case MSE $25.1707\to4.7808$, oracle MSE $3.4580$, and probability movement $0.03656\to0.07912$. The comparison depends on the metric: mean absolute off-target implied-log-odds movement decreases from $3.0506$ to $2.8866$ in the original and from $3.0361$ to $2.9091$ on CPU. The claim is increased probability movement, not an increase in every off-target metric. The artifact retains both records and the CPU invocation under \path{hmm/frozen/specificity_cpu_verification/}; the original invocation is unavailable, so the CPU check verifies the released defaults rather than an identical historical random stream.

The trained-readout and exact-posterior target MSEs are $3.441$ and $3.484$ in the original fixed-direction sweep, compared with $3.4862$ and $3.4580$ in the CPU verification. The small ordering reverses, so the records support comparable control performance, not a stable advantage for the trained readout.

\subsection{HMM recovery and calibration designs}\label{app:hmm-measurements}

Both studies define 32 components by crossing four layers with eight contiguous position bins. For the 63 input positions, the bin boundaries are $0,7,15,23,31,39,47,55,63$. Each layer has a unit high-minus-low $z_1$ direction, estimated from 20 batches of 256 sequences with at most 120,000 sampled positions and the same 20\% quantiles as above. A mask coefficient changes every position in that bin after the corresponding block. Because later activations depend on earlier edits, the coefficient-vector norm need not equal the norm of the total activation change through the network.

\paragraph{Aggregate recovery.}
The retained \path{hmm/frozen/nt_mi_set1_v2/mt_summary.json} supplies the 256 masks used in Sec.~\ref{sec:aggregate}. Thirty-two positive singleton interventions provide the reference. Responses are not divided by $0.6$, so fitted coefficients describe effects at that intervention size.

The sample-efficiency analysis permutes the saved rows with NumPy seed 7, using the validation and test splits in Sec.~\ref{sec:aggregate} and a nested fitting pool of 128. Fitting budgets are 8, 12, 16, 24, 32, 48, 64, 96, and 128. OMP selects support size from $\{1,2,3,4,5,6,8,10,12,16\}$; ridge uses penalty $0.01$, and Lasso selects among 24 penalties on the same validation set. Fits include an intercept. Pearson correlation compares 32 recovered effects with the singleton reference. The retained table, \path{hmm/frozen/nt_mi_sparse_v1/sparse_recovery_sample_efficiency.csv}, records 76 non-test measurements for the 12-fitting-mask OMP result, versus 96 for the 32-fitting-mask ridge comparison. Acquiring the reference and scoring tests are additional costs.

\paragraph{Calibration.}
Directions use seed 7 throughout the crossed experiment in Sec.~\ref{sec:step0}. Separate random streams generate masks and their split. Its unit-direction scale is not a fraction of the activation norm.

Let $v=\eps A\bar x$ be the attribution prediction and $y$ the measured, unnormalized response. The gain fit is the no-intercept least-squares correction $\widehat g=v^\top y/(v^\top v)$, with the code's small denominator guard. OMP selects support size from 1 through 12 using the 24 validation masks. The small-gain-budget comparison draws 4, 8, or 16 masks from the 96 non-test masks, repeats each draw 50 times using base seed 1701, and also fits on all 96. Test $R^2$ is computed over the 32 batch means, not individual sequences. Leave-one-batch-out calibration fits on the two other context batches at a fixed mask design. The saved configuration and summary are under \path{revision_checks/frozen/calibration_seed_factorial_disjoint/}; the earlier selector replay uses its own retained pools and fixed schedule, described in Appendix~\ref{app:selector-replay}.

\subsection{Synthetic response, masks, noise, and HVP queries}\label{app:synthetic-setup}

Equation~\ref{eq:synthetic-response} gives the retained response function using one-based indices; the code uses indices 0 through 6 for the active components and distractor. Its off-diagonal pair curvature is constant, although its saturating main effects make the full response nonquadratic. An additive-plus-pair fit approximates these main effects, leaving their residual in $w$.

Mask generation follows Sec.~\ref{sec:interactions}, replacing an all-zero row with one signed singleton before normalization. The correlated-design test copies coordinate 1 into coordinate 7 with probability $0.85$ before normalization; test rows never force this copying. Noise is added independently to fitting and test responses.

The model-class comparison uses 256 measurements at scales $0.6,1.2,2,5,8$. The budget sweep at scales 5 and 8 uses $16,24,32,48,56,64,72,80,96,112,128,192,256$ measurements, including validation as specified in Sec.~\ref{sec:interactions}. First-order OMP considers up to 12 terms and lifted OMP up to 16 among all 2,080 features, using $10^{-8}$ regularization in support fits; dense ridge and multi-gain fits use $0.001$. Singleton effects $\beta_i\tanh(\lambda_i\eps)/(\lambda_i\eps)$ are exact coordinate responses but omit pair effects and changes in a main effect's slope under differently weighted masks.

Prediction is test $R^2$ over normalized responses; pair recall is the fraction of the two true pairs present among the two largest-magnitude fitted pair coefficients. The 64- and 72-measurement thresholds apply the stated $R^2\ge0.95$ and recall $\ge0.99$ criteria to means over the three seeds. At a fixed budget of 96, the noise sweep uses $\sigma\in\{0.01,0.05,0.1,0.2,0.3\}$. Saved argument dictionaries and results are in the corresponding \path{planted_interactions/frozen/claim3_*/} directories.

The white-box comparison simulates noisy HVP access analytically. Its matrix is the Hessian of $\widetilde y_{\eps}$ with respect to mask weights at zero: off-diagonal entries are $\eps$ times the two planted pair coefficients, and diagonal entries are zero because $\tanh''(0)=0$. Each unit-norm signed query vector, with activation probability $0.3$, returns the full 64-vector $Hv$ plus independent Gaussian noise of standard deviation $0.01$ per entry. Query budgets are $1,2,4,8,12,16,24,32,48,64$; 25\% of queries select an OMP support of at most 16 pairs, with ridge penalty $10^{-8}$. The deterministic random seed is $s+1000\eps+17m$ for seed $s\in\{0,1,2\}$ and query budget $m$. Evaluation uses 512 new forward masks with response noise $0.01$. The matrix construction and noiseless gradient are simulated access primitives, not timed automatic-differentiation calls through a transformer. Each HVP returns 64 scalars, so its query count is not a forward-pass-equivalent cost. The retained configuration is \path{hvp_interactions/frozen/claim3_hvp_sparse030/claim3_hvp_metadata.json}.

\subsection{Qwen configuration and statistical reproduction}\label{app:qwen-reproduction}

The retained full configuration uses Qwen-2.5-7B-Instruct with the pinned revision
\begin{quote}
\path{a09a35458c702b33eeacc393d103063234e8bc28}.
\end{quote}
Loading uses bfloat16, batch size 32, and maximum sequence length 384. Directions at the eight layers in Sec.~\ref{sec:qwen} have norm $5\%$ of the construction activation scale. Design seed 2501 generates 192 fitting, 64 validation, and 128 test actions, with densities $0.25,0.5,0.75,1.0$; clean and singleton controls complete the 401-action design. Fitting and validation use scales $0.5$ and $1.0$, while testing uses $0.75$.

The four refusal stems are ``I cannot assist with that.'', ``I'm sorry, but I can't help with that.'', ``I can't provide instructions for that.'', and ``I must refuse this request.'' The compliance stems are ``Sure, here is the information you requested:'', ``Certainly. Here are the steps:'', ``Of course. First,'', and ``Here is how to do it:''. Both maps choose their ridge penalty from $\{10^{-6},10^{-4},0.01,0.1,1,10,100\}$. The paired bootstrap uses 5,000 resamples, seed 2529, significance level $0.05$, and the predeclared $5\%$ practical-improvement threshold. The prepared prompt manifest, frozen measurement records, and CPU analysis entry point are documented in \path{qwen/docs/REPRODUCING.md}; the precision audit is in Appendix~\ref{app:precision}.
